\documentclass{ctexart}

% 如果 XeTeX/LuaTeX 符号未自动定义，保留此定义
\providecommand*{\XeTeX}{X\kern-.1667em\lower.5ex\hbox{E}\kern-.125em\TeX}
\providecommand*{\LuaTeX}{Lua\TeX}

\begin{document}
当 \TeX 处理输入内容时，它不仅会读取单个字符，还会为每个字符分配一个类别。这个类别决定了该字符应如何被处理。这种分配是通过所谓的 “类别代码”（简称 “catcodes”）来完成的。每个输入字符都对应一个字符代码，并且每个字符代码都被分配给一个（可更改的）类别代码。

总的来说，\TeX 识别16种不同的类别。在文档中，一个字符所属的类别可以发生变化。最常见的例子是与语言相关的行为，比如由babel[1]构建的内容。在德语文档中，输入"a会生成“ä”。（在英语文档中，每个字符是单独处理的。）这是通过第13类“活动字符（active）”来实现的。活动字符不再是简单的字符，而是命令；在这种情况下，它是为后面的字符添加变音符号的命令。

\LaTeX 使用@来保护内部宏，防止终端用户访问。通常，内部宏以这种方式受到保护是有原因的。因此，在进行自定义操作时应谨慎，以避免意外的副作用。虽然要意识到其中的风险，但通过将@的类别代码更改为11（字母），并在不再需要时改回12（其他字符），就可以在命令名称中使用@： 
\end{document}